\documentclass{ximera}

\title{Trigonometric Derivatives Activity Facilitator Guide}
\author{MATH 425: Calculus I}

\begin{document}
\begin{abstract}
    Working with the peers in your group, solve the following problems. Make sure to show and justify all your work. Make sure everyone in the group understands the solution and participates. Be prepared to report your answers to the whole class. 
\end{abstract}
\maketitle

\textbf{Student Learning Outcomes}
Students will be able to: 
    \begin{itemize}
        \item Verify the derivatives of the reciprocal trig functions using the reciprocal defintions.
        \item Calculuate derivatives (and second derivatives) of functions involving trig functions.
        \item Interpret the meaning of the derivative based on context.
        \item Write the equation of a tangent line to a function involving trig functions.
    \end{itemize}

\textbf{Prerequisite Knowledge:} This activity does not explicitly require use of the chain rule, but certainly provides options for students who choose to use it (different books cover this rule in different places).  Students should be familiar with the quotient rule and the derivatives of sine, cosine, and tangent.  Students will also need the reciprocal definitions of cosecant, secant, and cotangent.\\

This worksheet may be long for some students.  Please note that question one only intends students to try one of three, but groups that are very fast can go back and finish off the other two.  The last two problems are relatively straightforward, and can be used as additional practice by students who do not finish in class.  

\begin{exercise}
    Confirm the derivatives of the following trigonometric functions using the quotient rule and the reciprocal definitions of these trig functions.  \textbf{Choose one of the three} to try! 
    \begin{enumerate}
        \item $\frac{d}{dx}(\sec(x))=\sec(x)\tan(x)$\\
        \textcolor{blue}{ $\frac{d}{dx}(\sec(x))=\frac{d}{dx}(\frac{1}{\cos(x)})$\\
        Using the quotient rule: $=\frac{(\cos(x)(0)-(-\sin(x))}{\cos^2(x)}$\\
        $=\frac{\sin(x)}{\cos(x)\cos(x)}=\frac{1}{\cos(x)}\frac{\sin(x)}{\cos(x)}=\sec(x)\tan(x)$\\
        Using chain rule and negative exponent: $=\frac{d}{dx}((\cos(x))^{-1})$\\
        $=-(\cos(x))^{-2}(-\sin(x))=\frac{\sin(x)}{\cos^2(x)}$ (see above for full simplification)}

        \item $\frac{d}{dx}(\csc(x))=-\csc(x)\cot(x)$\\
        \textcolor{blue}{ $\frac{d}{dx}(\csc(x))=\frac{d}{dx}(\frac{1}{\sin(x)})$\\
        Using the quotient rule: $=\frac{(\sin(x)(0)-(\cos(x))}{\sin^2(x)}$\\
        $=\frac{-\cos(x)}{\sin(x)\sin(x)}=-\frac{1}{\sin(x)}\frac{\cos(x)}{\sin(x)}=-\csc(x)\cot(x)$\\
        Using chain rule and negative exponent: $=\frac{d}{dx}((\sin(x))^{-1})$\\
        $=-(\sin(x))^{-2}(\cos(x))=-\frac{\cos(x)}{\sin^2(x)}$ (see above for full simplification)}

        \item $\frac{d}{dx}(\cot(x))=-\csc^2(x)$\\
        \textcolor{blue}{ $\frac{d}{dx}(\cot(x))=\frac{d}{dx}(\frac{\cos(x)}{\sin(x)})$\\
        Using the quotient rule: $=\frac{(\sin(x)(-\sin(x))-\cos(x)\cos(x)}{\sin^2(x)}$\\
        $=\frac{-\sin^2(x)-\cos^2(x)}{\sin^2(x)}=\frac{-(\sin^2(x)+\cos^2(x))}{\sin^2(x)}=\frac{-1}{\sin^2(x)}=-\csc^2(x)$\\
        Using chain rule and negative exponent: $=\frac{d}{dx}(\cos(x)(\sin(x))^{-1})$\\
        $=\cos(x)(-(\sin(x))^{-2}\cos(x))+(-\sin(x))(\sin(x))^{-1}$\\
        $=-\frac{\cos^2(x)}{\sin(x)^2}-1=\frac{-\cos^2(x)-\sin^2(x)}{\sin^2(x)}$(see above for full simplification)}
    \end{enumerate}
    \textcolor{orange}{Students do not necessarily need to fully simplify these derivatives in general, but we wanted to make the connection between the derivatives given in a textbook and the derivatives they may find from the rules.}
\end{exercise}

    
\begin{exercise}
    Find the derivative.
        \begin{enumerate}
            \item $g(x)=\frac{2x^4-3x^{-3}}{\sqrt[3]{x}}$\\
            \textcolor{blue}{ $g'(x)=\frac{(\sqrt[3]{x}(8x^3+9x^{-4})-(2x^4-3x^{-3})(\frac{1}{3}x^{-\frac{2}{3}})}{(\sqrt[3]{x})^2}$}

            \item $f(z)=(z^2-3z^{\frac{1}{4}})(z-1)^{-2}$\\
            \textcolor{blue}{ $=\frac{z^2-3z^{\frac{1}{4}}}{z^2-2z+1};\;\; f'(z)=\frac{(z^2-2z+1)(2z-\frac{3}{4}z^{-\frac{3}{4}})-(z^2-3z^{\frac{1}{4}})(2z-2)}{(z^2-2z+1)^2}$}\\
            \textcolor{orange}{This problem can be tackled with the chain rule if students already know it.  This solution assumes students do not yet know the chain rule.}

            \item $h(t)=\frac{e^t(3t-2)}{t^{\frac{3}{4}}}$\\
            \textcolor{blue}{ $h'(t)=\frac{t^{\frac{3}{4}}(3e^t+e^t(3t-2))-e^t(3t-2)(\frac{3}{4}t^{-\frac{1}{4}})}{t^{\frac{3}{2}}}$}

            \item $g(x)=\frac{x^7-3x^2+7x-4}{15}$ \\
            \textcolor{blue}{ $g(w)=\frac{7w^6-6w}{15}$}

            \item $f(t)=\frac{te^{2t}}{e^{3t}}$\\
            \textcolor{blue}{ $=\frac{t}{e^t};\;\; f'(t)=\frac{e^t-te^t}{e^{2t}}$}

            \item $p(x)=\frac{\sqrt{x}e^x}{x-3}$\\
            \textcolor{blue}{ $p'(x)=\frac{(x-3)(\sqrt{x}e^x+e^x(\frac{1}{2}x^{-\frac{1}{2}}))-\sqrt{x}e^x}{(x-3)^2}$}

            \item $f(\theta)=\theta\cos(\theta)$\\
            \textcolor{blue}{$f'(\theta)=\theta \sin(\theta)+\cos(\theta)$}

            \item $g(x)=\frac{\cos(x)-3x}{\sin(x)}$\\
            \textcolor{blue}{$g'(x)=\frac{\sin(x)(-\sin(x)-3)-(\cos(x)-3x)\cos(x)}{\sin^2(x)}$}

            \item $h(t)=\frac{\cos(t)}{2e^t}$ \\
            \textcolor{blue}{ $h'(t)=\frac{-2e^t\sin(t)-2e^t\cos(t)}{(2e^t)^2}$}

            \item $p(x)=2\sin(x)-\cos(x)$\\
            \textcolor{blue}{ $p'(x)=2\cos(x)+\sin(x)$}

            \item $g(\theta)=\theta e^\theta \sin(\theta)$\\
            \textcolor{blue}{ $g'(\theta)=\theta(e^\theta \cos(\theta)+e^\theta \sin(\theta))+e^\theta\sin(\theta)$\\}
            \textcolor{orange}{There is not a true need for a triple chain rule, but groups that are interested can try to derive one themselves (use function names and apply the product rule twice). }
    \end{enumerate}
\end{exercise}

\begin{exercise}
    When a mass hangs from a spring and is set in motion, the object's position oscillates in a way that the size of the oscillations decrease. This is usually called a \textit{damped oscillation}. Suppose that for a particular object, its displacement from equilibrium (where the object sits at rest) is modeled by the function \[s(t)=\frac{15\sin(t)}{e^t}.\]
    Assume that $s(t)$ is measured in inches and $t$ is measured in seconds. 
    \begin{enumerate} 
        \item Use Desmos to explore graph of this function for $t>0$ to see how it represents the situation described. \\
        \begin{center}
            \desmos{fqi3fp3vby}{800}{600}
        \end{center}
        \textcolor{blue}{The graph oscillates above and below the $t$-axis, which represents the equilibrium position. The size of the oscillations decrease as $t$ increases, which matches the description of a damped oscillation.}

        \item Compute $\frac{ds}{dt}$, state the units on this function, and explain what it tells you about the object's motion. \\
        \textcolor{blue}{$\frac{ds}{dt}=\frac{15e^t\cos(t)-15e^t\sin(t)}{e^{2t}}$\\
        This derivative tells the rate of change of the displacement, i.e. the velocity of the object at time $t$.}

        \item Compute and interpret $s'(2)$.\\
        \textcolor{blue}{$s'(2)=\frac{15e^2\cos(2)-15e^2\sin(2)}{e^{4}}=-2.69$ \\
        This means that the object is moving downward at a rate of 2.69 inches per second.}
    \end{enumerate}
    Problem adapted from Boelkins, et al., \textit{Active Calculus} 2nd ed.
\end{exercise}

\begin{exercise}
    Suppose that $f(x)=\frac{3\sin(x)}{e^x}$.  Find the equation of the tangent line at $x=0.$\\
    \textcolor{blue}{ $f(0)=\frac{3\sin(0)}{e^0}=\frac{3(0)}{1}=0$\\
    $f'(x)=\frac{3e^x\cos(x)-3\sin(x)e^x}{e^{2x}}$\\
    $f'(0)=\frac{3e^0\cos(0)-3\sin(0)e^0}{e^0}=\frac{3(1)(1)-3(0)(1)}{1}=3$\\
    Tangent line: $y-0=3(x-0)$ or $y=3x$}
\end{exercise}

\begin{exercise}
    Suppose that $g(x)=e^x(\sin(x))$.  Calculate $g'(x)$ and $g''(x)$. \\
    \textcolor{blue}{ $g'(x)=e^x\cos(x)+e^x\sin(x)$ (there are different options for how to simplify this - or not simplify - before using the product rule.)\\
    $g''(x)=e^x(-\sin(x))+e^x\cos(x)+e^x\cos(x)+e^x\sin(x)=2e^x\cos(x)$}
\end{exercise}  
   

\begin{exercise}
    Challenge: Differentiate $y=\cos^2(x)+\sin^2(x)$ in two different ways.  Confirm the answers match. \textcolor{blue}{ $0$ Hint: What is the Pythagorean identity? }\\
    \textcolor{orange}{Since the chain rule is not required on this worksheet, the first option would be to use the product rule on both $\cos^2(x)=\cos(x)\cos(x)$ and $\sin^2(x)=\sin(x)\sin(x)$.  The second option if students do not know the chain rule is the Pythagorean identity $\sin^2(x)+\cos^2(x)=1$.  Of course, the chain rule is available as an additional option if desired.}
\end{exercise}




\end{document}